% Generated draft: requires mathematical and source review.
\documentclass[11pt]{article}
\usepackage[margin=25mm]{geometry}
\usepackage{fontspec}
\setmainfont{Latin Modern Roman}
\usepackage{amsmath,amssymb,mathtools,unicode-math}
\setmathfont{Latin Modern Math}
\usepackage{enumitem,needspace,xurl,hyperref}
\setlength{\parindent}{0pt}
\setlength{\parskip}{5pt}
\setlength{\emergencystretch}{3em}
\newcommand{\N}{\mathbb{N}}
\newcommand{\Z}{\mathbb{Z}}
\newcommand{\Q}{\mathbb{Q}}
\newcommand{\R}{\mathbb{R}}
\newcommand{\C}{\mathbb{C}}
\newcommand{\F}{\mathbb{F}}
\newcommand{\A}{\mathbb{A}}
\newcommand{\PP}{\mathbb{P}}
\newcommand{\E}{\mathbb{E}}
\newcommand{\eps}{\varepsilon}
\newcommand{\dd}{\,\mathrm{d}}
\begin{document}
\section*{UnsolvedMath ID: 49}
\textbf{Source problem code:} catalogue entry
\subsection{Definitions and mathematical statement}
A height bound for rational points requires a specified ground field and height. For example, for a chosen projective embedding \(C\subset\PP^r_{\Q}\), a point \(P=[x_0:\cdots:x_r]\in C(\Q)\) has logarithmic projective height
\[
h(P)=\log\max_{0\leq i\leq r}|x_i|,
\]
when the \(x_i\) are relatively prime integer homogeneous coordinates. This example fixes neither the ground field nor the embedding for the original question; both are missing from its wording.

The supplied problem attributed to Stephen Smale (1998) asks for effective uniform bounds for heights of rational points on algebraic curves. In a chosen class of encoded curves and for a specified list \(I(C)\) of permitted invariants or presentation data, a precise version would require a computable function \(B\) satisfying
\[
\forall C\text{ in the chosen class}\ \forall P\in C(\Q),\qquad
h(P)\leq B(I(C)).
\]
Here effective means that the bound can be obtained by an algorithm from the permitted finite input, and uniform means that the same function works over the stated class with only the declared input dependencies. The class of curves, height, input invariants and dependencies are not fixed by the source, so the displayed implication is a framework for the intended problem rather than a completed replacement theorem.

The background invokes Faltings's finiteness result for curves of genus greater than one, suggesting the smooth projective high-genus setting; the displayed source question itself says arbitrary algebraic curves and does not impose that restriction. A finite bound for all rational points of an unrestricted curve cannot be assumed when the selected height is unbounded on those points. The source also moves between integer solutions in its motivation and rational points in its actual question; these are not silently identified.

The supplied literature triage dated 2026-08-17 labels the problem open, reports effective methods in specified classes or under additional hypotheses, and reports no verified uniform effective theorem in the requested generality. This is a dated source assessment rather than a newly verified current literature claim.
\subsection{Short English statement}
Make finiteness of rational points on suitable high-genus curves effective by giving a computable height bound with explicit uniform dependencies. The imported statement must specify the ground field, curve class, height, and permitted input data before that aim becomes one precise theorem.
\subsection{Sources}
\begin{itemize}
\item Local UnsolvedMath registry, permanent ID 49, source code catalogue entry, ``Smale's 5th Problem: Height Bounds for Diophantine Curves,'' attributed to Stephen Smale (1998), including supplied literature triage dated 2026-08-17.
\item Wolfram MathWorld, ``Smale's Problems,'' supplied access date 2026-08-17: \url{https://mathworld.wolfram.com/SmalesProblems.html}. The supplied triage cites this as a status list for the height-bounds task; no primary formulation or precise quantitative dependencies are supplied.
\end{itemize}
\end{document}
